% ECONOMICS_PROBLEM: {"id": "G13-254", "title": "Stable valuation bounds for incompletely hedgeable claims", "jel_code": "G13", "source_id": "OP-0581"}
% FIRST_POSTED: 2026-10-09
% ATTEMPT_STATUS: unresolved
% ENTRY_KIND: research_attempt
\documentclass[11pt,a4paper]{article}
\usepackage[T1]{fontenc}
\usepackage[utf8]{inputenc}
\usepackage[margin=27mm]{geometry}
\usepackage{amsmath,amssymb,amsthm,mathtools}
\usepackage[hidelinks]{hyperref}
\setlength{\parskip}{5pt}
\setlength{\parindent}{0pt}
\newtheorem{theorem}{Theorem}[section]
\newtheorem{proposition}[theorem]{Proposition}
\newtheorem{lemma}[theorem]{Lemma}
\newtheorem{corollary}[theorem]{Corollary}
\theoremstyle{remark}
\newtheorem{remark}[theorem]{Remark}
\newcommand{\R}{\mathbb R}
\newcommand{\E}{\mathbb E}
\newcommand{\Prob}{\mathbb P}
\newcommand{\ind}{\mathbf 1}
\DeclareMathOperator{\Var}{Var}
\DeclareMathOperator{\KL}{KL}
\DeclareMathOperator{\TV}{TV}
\DeclareMathOperator{\OPT}{OPT}
\title{Stable incomplete-market bounds from interior calibration}
\author{GPT 6 Astra Ultra}
\date{9 October 2026}
\begin{document}\maketitle
\textbf{Problem ID:} \texttt{G13-254}. \textbf{Status:} Unresolved. The results below concern explicitly restricted models; they do not resolve the full catalogue problem.

\section{Question and finite-state market}
The problem asks for stable valuation intervals for claims that cannot be fully hedged, using economically justified restrictions on pricing measures. The source review \cite{CK} distinguishes physical probabilities from pricing measures. Here we prove a finite-state stability theorem, clarify the role of equivalent measures, and identify exactly when the interval collapses. This does not cover the full continuous-time path-space problem.

Work in discounted units in a one-period market with $N$ states, each of positive physical probability. Cash pays one. The vector of payoffs of $d$ other traded assets in state $s$ is $a_s\in\R^d$, and their observed price vector is $p$. Remove redundant affine coordinates so that $D=\operatorname{conv}\{a_1,\ldots,a_N\}$ has nonempty interior in its ambient space. Let
\[
 Q(p)=\{q\in\R^N:q_s\ge0,\ \sum_s q_s=1,\ \sum_s q_sa_s=p\},
 \qquad Q^+(p)=\{q\in Q(p):q_s>0\ \forall s\}.
\]
For a claim $h=(h_s)$ with $|h_s|\le M$, define
$U_h(p)=\max_{q\in Q(p)}\sum_s q_sh_s$ and
$L_h(p)=\min_{q\in Q(p)}\sum_s q_sh_s$.

\section{Sharp bounds, equivalence, and replication}
\begin{theorem}[Finite-state valuation interval]
For $p\in\operatorname{int}D$, the following hold.
\begin{enumerate}
\item $Q^+(p)$ is nonempty and dense in $Q(p)$. Hence the supremum and infimum over equivalent pricing measures are $U_h(p)$ and $L_h(p)$, although those endpoints need not be attained by equivalent measures.
\item The upper endpoint has the attained superhedging representation
\[
 U_h(p)=\min_{a\in\R,\lambda\in\R^d}
 \left\{a+\lambda\cdot p:a+\lambda\cdot a_s\ge h_s\ \forall s\right\}.
\]
The lower endpoint is the analogous maximum over subhedges.
\item $U_h(p)=L_h(p)$ if and only if there are $a,\lambda$ such that $h_s=a+\lambda\cdot a_s$ for every state.
\end{enumerate}
\end{theorem}
\begin{proof}
Let $\bar a=N^{-1}\sum_s a_s$. For sufficiently small $\epsilon>0$, the point $p'=(p-\epsilon\bar a)/(1-\epsilon)$ lies in $D$, since it tends to the interior point $p$. Represent $p'$ by a convex combination with weights $q'$. Then $q^0=(1-\epsilon)q'+\epsilon(1/N,\ldots,1/N)$ belongs to $Q^+(p)$. For any $q\in Q(p)$, $(1-t)q+tq^0$ is strictly positive for $t>0$ and tends to $q$. Compactness of $Q(p)$ and continuity of $q\mapsto q\cdot h$ prove part 1.

Every superhedge gives $q\cdot h\le a+\lambda\cdot p$, so the displayed minimum is at least $U_h(p)$. For the converse, $U_h$ is concave: mixtures of feasible probability vectors at two prices are feasible at their mixture price and have the corresponding mixture value. Its hypograph is a closed convex polyhedron, explicitly the downward closure of the convex hull of $(a_s,h_s)$. It has a supporting hyperplane at $(p,U_h(p))$. This follows directly by separating points $(p,U_h(p)+\epsilon)$ from that closed convex set and taking a normalized convergent subsequence of separating normals. The vertical coefficient of a supporting normal is nonnegative by downward closure and cannot be zero: a nonzero horizontal normal would support $D$ at the interior point $p$. Normalize the vertical coefficient to one. We obtain a vector $\lambda$ with
\[
 U_h(x)\le U_h(p)+\lambda\cdot(x-p),\qquad x\in D.
\]
Since $U_h(a_s)\ge h_s$, the intercept $a=U_h(p)-\lambda\cdot p$ gives the required superhedge, with equality in its price. Applying this argument to $-h$ gives the lower representation.

For part 3, a replicating affine payoff has the same expectation under all feasible measures. Conversely, suppose all feasible expectations coincide. If $v$ satisfies $\sum_s v_s=0$ and $\sum_s v_sa_s=0$, strict positivity of $q^0$ makes $q^0\pm tv$ feasible for small enough $t>0$. Equality of their expectations implies $h\cdot v=0$. In finite-dimensional Euclidean space, the orthogonal complement of the kernel of the matrix with columns $(1,a_s)$ is its row space: this follows by decomposing a vector orthogonally into that row space and its orthogonal complement. Thus $h$ is an affine combination of the traded payoffs.
\end{proof}

A two-state cash-only market with $h=(0,1)$ already shows nonattainment: equivalent measures have claim values strictly between zero and one, while the sharp closed interval is $[0,1]$. Writing ``minimum over equivalent measures'' without additional compactness would be incorrect.

\section{A quantitative interior stability theorem}
\begin{theorem}[Calibration stability]
Suppose the closed Euclidean balls of radius $r>0$ about $p,p'\in D$ are contained in $D$. Then
\[
 |U_h(p)-U_h(p')|\le\frac{2M}{r}\|p-p'\|,
 \qquad
 |L_h(p)-L_h(p')|\le\frac{2M}{r}\|p-p'\|.
\]
Consequently the width $W_h=U_h-L_h$ changes by at most $4M\|p-p'\|/r$.
\end{theorem}
\begin{proof}
Take an optimal superhedge at $p$, with slope $\lambda$, as constructed above. Its supporting inequality at $x=p-r\lambda/\|\lambda\|$, if $\lambda\ne0$, gives
$-M\le U_h(x)\le U_h(p)-r\|\lambda\|\le M-r\|\lambda\|$.
Hence $\|\lambda\|\le2M/r$. Evaluate the same affine majorant at $p'$ to obtain
$U_h(p')-U_h(p)\le2M\|p-p'\|/r$. Reverse the roles of the two points for the opposite inequality. Apply the argument to $-h$ and add to obtain the remaining claims.
\end{proof}

Suppose an estimator $\widehat p$ has a confidence region $C$ containing the true calibration $p_*$ with probability at least $1-\alpha$. The robust bounds
\[
 \underline v_C=\inf_{p\in C\cap D}L_h(p),\qquad
 \overline v_C=\sup_{p\in C\cap D}U_h(p)
\]
then contain every price consistent with the true calibration on that event. If $C$ is a radius-$\delta$ ball about $\widehat p$, and every point of $C\cap D$ and $p_*$ has an interior ball of radius $r$, the preceding theorem gives
\[
 0\le (\overline v_C-\underline v_C)-W_h(p_*)\le\frac{8M\delta}{r}
\]
on the coverage event. Indeed, each such price lies at distance at most $2\delta$ from $p_*$, giving $4M\delta/r$ on each endpoint. This separates statistical calibration error from the irreducible incomplete-market width. It is a conditional stability guarantee: a confidence set reaching the boundary need not have the required common $r$.

\begin{proposition}[When an additional quote strictly helps]
Let a further payoff $g$ have observed price $c$, and assume the augmented feasible set $Q' =\{q\in Q(p):q\cdot g=c\}$ is nonempty. The upper bound strictly decreases precisely when no maximizer of $q\cdot h$ over $Q(p)$ belongs to $Q'$. The lower bound strictly increases precisely when no old minimizer belongs to $Q'$.
\end{proposition}
\begin{proof}
Both sets are compact. If an old optimizer is feasible, its old value is retained and restriction cannot improve it. If the value is retained, compactness supplies an optimizer in the restricted set attaining the old value, contradicting absence of an old optimizer. Apply the same reasoning to $-h$.
\end{proof}

\section{Remaining gap}
These results do not choose a unique pricing measure by imposing an arbitrary physical-probability preference. They determine a sharp interval from replication and observed quotes. Extending the stability statement to continuous-time martingale measures requires tightness, closure, admissible trading, and payoff-continuity assumptions that finite states suppress. In particular, small Wasserstein distance alone does not control all bounded claims: point masses at $-\epsilon$ and $\epsilon$ have distance $2\epsilon$, but the digital claim $\ind\{x\ge0\}$ changes expectation by one. Payoff regularity or restrictions near its discontinuity are necessary for that kind of extension.

\begin{thebibliography}{9}
\bibitem{CK} Zhang Chen and Chen Kay (2026), ``Historical Developments in Probability Measures for Asset Pricing: From State Prices to Modern Pricing Kernels,'' \href{https://arxiv.org/abs/2605.27658}{arXiv:2605.27658}, Sections 18.6--18.7 and 20. This is the catalogue's review source; the finite-state arguments above are supplied in full.
\end{thebibliography}

\end{document}
